\section{Audit Chain and Legal Properties}
\label{sec:audit}

\subsection{The Audit Chain Record}

Every \vr{} run appends a structured record to
\texttt{runs/\{label\}/\{year\}/index.json}:

\begin{verbatim}
{
  "search": "vra-recom",
  "forest_steps": 1000,
  "percentile": 0.0,
  "base_seed": 12345678,
  "vap_threshold": 0.50,
  "protected_districts": [1, 12],
  "steps_taken": 1000,
  "steps_accepted": 412,
  "vra_rejections": 231,
  "vra_rejection_rate": 0.23,
  "mh_rejections": 357,
  "selected_step": 291,
  "forward_seed_formula":
    "SHA-256('FR_FORWARD_'||step:u64le||'_'||0u32le||'_'||base_seed:u64le)",
  "reverse_seed_formula":
    "SHA-256('FR_REVERSE_'||step:u64le||'_'||0u32le||'_'||base_seed:u64le)"
}
\end{verbatim}

The record contains all information needed to independently reproduce the run:
the base seed, the forward/reverse seed derivation formulas, the protected
district list, and the exact step counts.
The accounting identity
$\text{steps\_accepted} + \text{vra\_rejections} + \text{mh\_rejections} =
\text{steps\_taken}$ is verifiable from the record without re-running the chain.

\subsection{Auditor Verification Protocol}

An auditor---a court-appointed expert, opposing counsel, or an independent
researcher---can verify any \vr{} run using only the public audit record:

\begin{enumerate}
  \item \textbf{Seed derivation.}
        Re-derive all forward and reverse seeds from \texttt{base\_seed} using
        the formula in the record.
        These are deterministic functions of \texttt{base\_seed} and the step
        index; no other state is required.

  \item \textbf{Proposal replay.}
        Run a \fr{} chain (no VRA checking) with the same \texttt{base\_seed}.
        Because \vr{} inherits \fr{}'s seed derivation, the proposal at each
        step is identical in both chains.
        The auditor can identify exactly which steps \vr{} accepted, which
        it rejected via MH, and which it rejected via VRA.

  \item \textbf{VRA verification.}
        For each step accepted by \vr{}, verify that all districts in
        \texttt{protected\_districts} have minority VAP $\geq$
        \texttt{vap\_threshold} under the accepted plan.
        This is a straightforward computation from the accepted plan
        assignment and the public census VAP data.

  \item \textbf{Step count verification.}
        Confirm that the counts in the record satisfy the accounting identity.
        Confirm that the plan at \texttt{selected\_step} (the $p$-th percentile
        of the accepted plans by EC) matches the submitted plan.
\end{enumerate}

This protocol does not require access to the redist binary or any proprietary
code: only the census data, the SHA-256 implementation, and arithmetic are
needed.
The auditability property follows directly from the seed inheritance and the
determinism of the algorithm.

\subsection{Legal Properties}

\paragraph{Conditional distribution correctness.}
The target distribution of \vr{} is the conditional distribution over
$\Pi_{\mathrm{vra}}(G, k, \Prot)$: valid plans in which every protected
district has minority VAP $\geq \theta$.
By Proposition~\ref{prop:vra-invariant}, VRA compliance is a chain invariant.
By the \fr{} expected detailed balance theorem~\citep{dellaLibera2026forestRecom},
the chain's stationary distribution on $\Pi_{\mathrm{vra}}$ is the
conditional uniform distribution restricted to the VRA-compliant sub-space,
in the expected-detailed-balance sense.

\paragraph{Conditional distributional correctness.}
Proposition~\ref{prop:vra-invariant} establishes that the VRA check preserves the $\Pi_{\mathrm{vra}}$ invariant. We now show this does not distort the expected detailed balance property of \fr{} on $\Pi_{\mathrm{vra}}$. Let $\pi, \pi' \in \Pi_{\mathrm{vra}}$ be two VRA-compliant plans. The VRA-Aware chain proposes the move $\pi \to \pi'$ with the same probability as unconstrained \fr{}: $Q_{\mathrm{vra}}(\pi, \pi') = Q_{\mathrm{FR}}(\pi, \pi')$. It accepts with probability $\min(1, c_{\mathrm{fwd}}/c_{\mathrm{rev}})$ when VRA-compliant, and 0 otherwise. For moves within $\Pi_{\mathrm{vra}}$, the acceptance probability and proposal probability are identical to constrained FR, so the expected detailed balance with respect to the uniform distribution on $\Pi_{\mathrm{vra}}$ follows from the FR result (\citet{mccartan2020} \S3.2). The VRA hard-rejection on moves leaving $\Pi_{\mathrm{vra}}$ reduces the effective step rate but does not introduce bias on compliant plans.

\paragraph{Reference plan matching.}
The protected set $\Prot$ is derived from the \emph{initial plan} $\pi_0$,
not from the current chain state.
This matches the legal standard: the obligation is to preserve the
majority-minority districts that appear in the submitted plan.
Courts and expert witnesses can confirm that \texttt{protected\_districts}
in the audit record matches the district IDs of the majority-minority districts
in the submitted plan, closing the loop between the algorithmic protection
and the legal requirement.

\paragraph{Relationship to VRASection.}
\vraSection{} and \vr{} are both necessary for a complete VRA-compliant
litigation analysis:
\begin{itemize}
  \item \vraSection{} generates the submitted plan: a valid plan with
        majority-minority districts.
  \item \vr{} generates the comparison ensemble: alternative valid plans
        that also preserve those districts.
\end{itemize}
Using only \vraSection{} without \vr{} provides a single plan but no ensemble
context.
Using \vr{} without \vraSection{} risks starting from a non-compliant plan,
causing an initial burn-in period before the chain reaches the compliant
sub-space (which, by Proposition~\ref{prop:vra-invariant}, it can never leave
once the initial plan is non-compliant---since the VRA check rejects transitions
that would \emph{reduce} VAP below threshold, but cannot enforce compliance
on plans already below threshold).
The recommended composite is:

\begin{verbatim}
structure: ratio-optimal-vra   # VRASection: compliant starting plan
weights: vra-aligned           # load minority VAP weights
search: vra-recom              # VraRecomChain: compliant ensemble
vra_threshold: 0.50
\end{verbatim}

\subsection{Warning: Empty Protected Set}

If no district in the initial plan has minority VAP $\geq \theta$, the CLI
emits:

\begin{verbatim}
WARNING: VRA-aware chain constructed with 0 protected districts --
VRA enforcement is a no-op. Check --vra-threshold or use
--weights-override vra-aligned to produce an initial VRA-compliant plan.
\end{verbatim}

This warning is critical for legal use: a practitioner who inadvertently
supplies a non-VRA-aligned initial plan will receive an ensemble that is
indistinguishable from an unconstrained \fr{} run.
The warning is unconditional (not suppressed by \texttt{--quiet}) to prevent
silent misuse.
